En el capítulo anterior mencionamos los jueces en línea (online judges), a grandes razgos, estos son sistemas que reciben un programa enviado por un usuario, entonces ejecutan este con entradas predefinidas a las que se les llama casos y compara las respuestas arrojadas por el programa con las esperadas, y a partir de esto determina si la solución enviada es correcta o no.
\\
El funcionamiento de un juez en línea se puede ver de la siguiente manera:
\begin{center}
    Cliente $\rightarrow$ Frontend $\rightarrow$ Backed $\rightarrow$ Evaluador
\end{center} 